% cv_skeleton.tex — format-by-example for /write-cv
% This file is package-shipped and refreshed by `init --refresh`.
% Edit slot bodies below to guide the /write-cv skill.

%% SLOT: recipient_company
% Empfänger-Firma oder -Organisation (DIN 5008 Zeile 1).
% Fett gesetzt im Anschriftenfeld. Sollte aus der Stellenanzeige hervorgehen.
Muster GmbH

%% SLOT: recipient_name
% Ansprechperson (Vorname Nachname) oder Abteilung (DIN 5008 Zeile 2).
% Wenn kein Name bekannt: "Personalabteilung" oder ganz leer lassen
% (leerer Body wird im PDF ohne Lücke übersprungen).
Frau Dr. Müller

%% SLOT: recipient_street
% Straße und Hausnummer (DIN 5008 Zeile 3).
Musterstraße 42

%% SLOT: recipient_zip_city
% Postleitzahl und Stadt (DIN 5008 Zeile 4).
20095 Hamburg

%% SLOT: opening
% Förmliche Anrede passend zum recipient_name (bzw. recipient_company).
% Bei unbekanntem Empfänger: "Sehr geehrte Damen und Herren,"
Sehr geehrte Frau Dr. Müller,

%% SLOT: cover_intro
% 2–3 Sätze: Wer bin ich, welche Stelle interessiert mich, warum dieses Unternehmen?
% Direkte Ansprache, keine Floskeln. Stärksten Bezugspunkt zur Stelle als Aufhänger.
mit großem Interesse habe ich Ihre Stellenausschreibung als Senior Data Engineer gelesen.
Die Verbindung von skalierbaren Datenpipelines und \textit{maschinellem Lernen} entspricht
genau meinem Schwerpunkt der letzten fünf Jahre --- und Ihr Fokus auf öffentliche Dateninfrastruktur
spricht mich besonders an.

%% SLOT: cover_pivot
% 1–2 Sätze: Brücke vom bisherigen Werdegang zur neuen Stelle.
% Warum ist dieser Wechsel der logische nächste Schritt?
Nach meiner Tätigkeit bei der Muster GmbH, wo ich vollständige ETL-Strecken von der
Rohdatenerfassung bis zum Reporting verantwortet habe, möchte ich dieses Know-how
nun in einem stärker forschungsnahen Umfeld einsetzen.

%% SLOT: cover_fit
% 2–3 Sätze: Konkrete Belege für den Fit — Zahlen, Technologien, Erfolge.
% Bezug auf die Anforderungen der Stellenausschreibung herstellen.
Mein Open-Source-Beitrag \href{https://github.com/beispiel/pipeline-lib}{pipeline-lib}
wird von über 300 Teams produktiv genutzt und zeigt, dass ich nicht nur implementiere,
sondern auch dokumentiere und community-tauglich designt. Umlaute wie Ü, Ö, Ä und ß
sind in meinen deutschen Texten selbstverständlich korrekt gesetzt.

%% SLOT: cover_closing
% 1–2 Sätze: Gesprächswunsch äußern, Verfügbarkeit andeuten, freundlicher Abschluss.
Über die Möglichkeit, meine Erfahrungen in einem persönlichen Gespräch vorzustellen,
würde ich mich sehr freuen --- idealerweise ab dem 1. Juni.

%% SLOT: resume_berufserfahrung
% Listet relevante Berufsstationen in umgekehrter chronologischer Reihenfolge.
% Je Station ein \cventry. Schwerpunkt auf Aufgaben mit Stellenbezug.
\cventry{2021--heute}{Senior Data Engineer}{\textit{Muster GmbH}}{Hamburg}{}{%
  Aufbau einer Echtzeit-Datenpipeline für 50\,Mio. Ereignisse/Tag;
  Einführung von dbt und Apache Kafka --- Laufzeit um 40\,\% reduziert.}
\cventry{2018--2021}{Data Engineer}{Beispiel AG}{Berlin}{}{%
  Entwicklung von ETL-Prozessen in Python; Anbindung externer APIs über
  \href{https://beispiel.de/api}{REST-Schnittstellen}.}

%% SLOT: resume_ausbildung
% Akademische Abschlüsse und relevante Zertifikate, neuestes zuerst.
% Promotion, Master, Bachelor; Berufsausbildung nur wenn fachlich relevant.
\cventry{2015--2018}{M.\,Sc.\ Informatik}{Universität Hamburg}{Hamburg}{}{%
  Schwerpunkt: Verteilte Systeme und maschinelles Lernen;
  Abschlussnote \textit{sehr gut} (1,2).}
\cventry{2012--2015}{B.\,Sc.\ Wirtschaftsinformatik}{FH Lübeck}{Lübeck}{}{}

%% SLOT: resume_projekte
% Eigenprojekte, Open-Source-Beiträge oder Nebenarbeiten mit technischem Belang.
% Projekttitel verlinken wenn öffentlich zugänglich (\href).
\cventry{2023}{pipeline-lib}{Open Source}{}{}{%
  Python-Bibliothek für deklarative Datenpipelines ---
  \href{https://github.com/beispiel/pipeline-lib}{github.com/beispiel/pipeline-lib};
  300\,+ produktive Installationen, MIT-Lizenz.}

%% SLOT: skills_block
% Mechanisch zusammengesetzt aus user-info/search-terms/skills.md per ADR-0033.
% Nicht hier hand-authoren --- /write-cv emittiert die \cvitem-Zeilen anhand
% der Gruppen-Auswahlregeln (always-Gruppen + LLM-Picks nach Jobtype-Relevanz)
% und der Item-Auswahlregeln (always-Items + LLM-Picks innerhalb gewählter Gruppen).
% File-Reihenfolge in skills.md = Render-Reihenfolge.
